\section{恒定磁场}
\sourceproblem{选择1}{31}{37}
\problemtext 同电流的圆形线圈半径 $a_1$、正方形线圈边长 $a_2$，中心磁感应强度大小相同，求尺寸比。
\workings 圆环中心 $B_1=\mu_0 I/(2a_1)$；正方形每边距中心 $a_2/2$，两端张角均45度，四边叠加得 $B_2=2\sqrt2\mu_0 I/(\pi a_2)$。令两者相等得 $a_1:a_2=\pi:4\sqrt2$。
\crosscheck 两式均随尺寸增大而减小；保留原书选项对应的二校标记，不以模糊选项替代公式。
\sourceproblem{选择2}{31}{37}
\problemtext 图2直角折线载流导线，两条半无限直边至O的垂距均为 $a$，求O点磁感应强度大小。
\workings 每一半无限边的近端相对垂足张角45度、远端张角90度，故各贡献 $\mu_0 I(1+1/\sqrt2)/(4\pi a)$。按图电流向右再向下，两段在O处磁场均垂直纸面向内。因此
\[
B=\frac{\mu_0 I}{2\pi a}\left(1+\frac1{\sqrt2}\right).
\]
\crosscheck 采用毕奥--萨伐尔有限直线公式，不能把两段都当成从垂足开始的半无限直线。
\sourceproblem{选择3}{31}{37}
\problemtext 均匀磁场 $\bm B$ 中，以半径 $R$ 的圆为边界张任意曲面 $S$，求磁通量。
\workings 取圆面法向 $\bm n$ 与曲面定向相容，即两曲面沿公共边界诱导方向相同。由 $\nabla\cdot\bm B=0$，两面磁通相等，$\Phi_S=\bm B\cdot\bm n\,\pi R^2=\pi R^2B\cos\alpha$，其中 $\alpha$ 为 $\bm B$ 与 $\bm n$ 的夹角。
\crosscheck 反转曲面定向则通量变号；若原图法向取闭合面圆底的外法向，应使用负号。图3箭头与选项的对应保留二校标记。
